\chapter{实验路线}

第二册的实验围绕 Compute Systems Labs 展开。实验主线采用事件流分析引擎，而不是只做独立日志行计数。最小层是 \code{event_name} 计数，用来建立 reference、split、partial 和 manifest；高性能层加入 \code{event_time}、\code{ingest_time}、\code{session_id}、\code{seq} 和 \code{object_id}，用来训练顺序依赖、会话状态、滑动窗口、乱序缓冲和维表 join；分布式层再把 partition、shuffle、迟到数据、幂等提交和 checkpoint 连接起来。第一组实验是硬件与测量基础，包括 \cmd{lscpu} 拓扑记录、缓存层级观察、TLB 与缺页计数、perf stat、顺序归约和并行归约。第二组实验是数据布局，包括 AoS、SoA、false sharing、顺序访问、随机访问和预取效果。第三组实验是单机高性能计算，包括多累加器、SIMD 自动向量化、Roofline 分析、典型内核模式、session scan 和 sliding window aggregate。

第四组实验是多线程同步，包括 mutex counter、atomic counter、bounded queue、条件变量、信号量和线程亲和性。第五组实验是无锁结构，包括 SPSC ring buffer、CAS 循环、ABA 现象、版本标签和内存回收方案分析。第六组实验是并行算法，包括 reduce、scan、partition、thread-local histogram、session-local state、窗口状态合并和任务粒度比较。第七组实验是运行时，包括线程池、任务队列、工作窃取模型、异步 I/O、背压和水位线推进。第八组实验是分布式计算，包括本机多进程 RPC 模拟、超时重试、幂等请求、分片、复制、shuffle、迟到事件修正和窗口 checkpoint。

\begin{labbox}
每个实验必须记录环境、输入、命令、输出、正确性检查和结论。性能数据没有上下文就不能进入正文。若实验受手机、Termux 或虚拟化环境限制，也要明确记录限制。
\end{labbox}
